% !TeX root = ../main.tex
\section{Introduction}\label{sec:intro}

Single-instruction-multiple-data (SIMD) units let one instruction operate on a whole vector of
elements. On the Intel Core i9-10900F studied here the widest unit is AVX2 with FMA: a
\qty{256}{bit} register holds \num{8} float32 or \num{4} float64 lanes, so the \emph{ideal} speedup
over scalar code is $8\times$ or $4\times$. Loads and stores, loop tails, reduction dependencies
and, above all, memory bandwidth keep real kernels below that. We measure where the benefit
appears, how far it falls short of the lane count, and when it disappears:
\begin{enumerate}
  \item How much faster are auto-vectorized builds than scalar builds for each kernel, and how does
        that change from cache-resident to DRAM-resident working sets? (\cref{sec:baseline,sec:locality})
  \item What do alignment, loop tails, stride and gather access, and data type cost? (\cref{sec:alignment,sec:stride,sec:dtype})
  \item Where do these kernels sit in a roofline model of this machine? (\cref{sec:roofline})
  \item When does the expected SIMD benefit fail to appear? (\cref{sec:negative})
\end{enumerate}
